\documentclass[11pt]{article}

\usepackage[margin=1in]{geometry}
\usepackage{amsmath, amssymb, amsthm}
\usepackage{booktabs}
\usepackage{array}
\usepackage{xcolor}
\usepackage{hyperref}
\usepackage{enumitem}
\usepackage{microtype}

\title{JEPA World Model for Visual Cartpole Control:\\
Architecture, Losses, and Data Generation}
\author{}
\date{}

\newcommand{\R}{\mathbb{R}}
\newcommand{\E}{\mathbb{E}}
\newcommand{\loss}[1]{\mathcal{L}_{\text{#1}}}
\newcommand{\zstar}{z^{*}}

\begin{document}
\maketitle

\section{Overview}

The JEPA world model (\texttt{JEPAModel}) consists of three learned components:
an \emph{online encoder} mapping pixel observations to a low-dimensional latent
state, a \emph{target encoder} (an EMA copy of the online encoder) that produces
stable prediction targets, an \emph{action encoder} that lifts the scalar control
input into the latent action space, and an \emph{MLP predictor} that models the
one-step latent dynamics. Critically, \textbf{JEPA has no decoder} — it is
trained purely in latent space via prediction and regularization losses (the
companion \texttt{AEWorldModel} variant adds a symmetric CNN decoder and is
described separately).

\section{Model Architecture}

\subsection{Online Encoder --- \texttt{CNNEncoder}}

Input: $(B, 6, 64, 64)$ --- two consecutive RGB frames stacked channel-wise
(\texttt{frame\_stack=2}, so $3 \times 2 = 6$ input channels).

Six convolutional blocks, each of the form
$\text{Conv2d}(3{\times}3,\ \text{pad}{=}1) \to \text{MaxPool}(2) \to \text{BN} \to \text{ReLU}$.
The final block omits batch-norm and ReLU because the spatial resolution collapses
to $1{\times}1$ (BatchNorm is undefined for spatial size 1 with batch size 1).

\begin{table}[h]
\centering
\begin{tabular}{@{}cccc@{}}
\toprule
Block & $C_{\text{in}} \to C_{\text{out}}$ & Spatial output & Notes \\
\midrule
0 & $6 \to 8$   & $32\times32$ & BN + ReLU \\
1 & $8 \to 16$  & $16\times16$ & BN + ReLU \\
2 & $16 \to 32$ & $8\times8$   & BN + ReLU \\
3 & $32 \to 64$ & $4\times4$   & BN + ReLU \\
4 & $64 \to 64$ & $2\times2$   & BN + ReLU \\
5 & $64 \to 8$  & $1\times1$   & \emph{no} BN, \emph{no} ReLU \\
\bottomrule
\end{tabular}
\caption{CNN encoder architecture. Output is flattened to $z \in \R^{8}$
(latent\_dim $= 8$). Approx.\ 17k parameters.}
\end{table}

\subsection{Target Encoder}

An exponential-moving-average (EMA) copy of the online encoder, updated every
training step:
\begin{equation}
\theta_{\text{target}} \leftarrow 0.99\,\theta_{\text{target}} + 0.01\,\theta_{\text{online}}.
\end{equation}
It produces stop-gradient targets $z_1, \ldots, z_H$ for the prediction loss
(computed under \texttt{torch.no\_grad()}); only $z_0$ is produced by the
\emph{online} encoder with gradients enabled. This is the standard
BYOL/I-JEPA mechanism for preventing representational collapse.

\subsection{Action Encoder --- \texttt{LinearActionEncoder}}

\begin{equation}
c_t = W u_t, \qquad W \in \R^{1 \times 1},
\end{equation}
i.e.\ a learned scalar multiplier on the scalar force $u_t$ (no bias,
\texttt{action\_latent\_dim=1}, \texttt{action\_encoder: linear}).

\subsection{Predictor --- \texttt{MLPPredictor} (window $W=3$)}

The predictor consumes the last $W=3$ latent states and encoded actions,
flattened into a single vector:
\begin{equation}
\text{input} = \big[\, z_{t-2},\, z_{t-1},\, z_t,\; c_{t-2},\, c_{t-1},\, c_t \,\big]
\in \R^{3 \times 8 + 3 \times 1} = \R^{27}.
\end{equation}
Architecture: $\text{Linear}(27 \to 64) \to \text{ELU} \to \text{Linear}(64 \to 8)$,
i.e.\ 2 layers with hidden dimension 64. Output: $\hat z_{t+1} \in \R^{8}$.

\bigskip
\noindent\textbf{Note:} \texttt{JEPAModel} contains \emph{no decoder}.
The full model is purely encoder $+$ predictor; the \texttt{AEWorldModel}
subclass adds a symmetric \texttt{CNNDecoder} (six transposed-convolution
blocks, channel progression $8 \to 64 \to 64 \to 32 \to 16 \to 8 \to 3$,
final Sigmoid activation) for pixel-space reconstruction and prediction
losses, but is a separate training pipeline (\texttt{train\_ae.py}).

\section{Training Data Generation}

The dataset is generated once and cached as HDF5. Episodes run to completion
with \emph{no} mid-episode resets --- the gym \texttt{done} signal is ignored
so trajectories show the pole falling freely past the $12^\circ$ threshold.

\subsection{Episode types (\texttt{cartpole\_v2\_fullspec.yaml})}

\begin{table}[h]
\centering
\small
\begin{tabular}{@{}lccccc@{}}
\toprule
Type & Episodes $\times$ Steps & Transitions & Init.\ range & Action policy \\
\midrule
Random   & $50 \times 200$  & 10{,}000 & $\pm 1.0$ rad ($\pm 57^\circ$) & $u \sim \mathcal{U}[-10, 10]$ \\
LQR      & $100 \times 100$ & 10{,}000 & $\pm 0.10$ rad ($\pm 17^\circ$) & GT-LQR, no noise \\
Eq-LQR   & $10 \times 50$   & 500      & $\pm 0.002$ rad ($\approx 0^\circ$) & GT-LQR $+\ \sigma{=}0.001$ \\
PRBS     & $125 \times 40$  & 5{,}000  & $\pm 0.05$ rad & $\pm 3$N, flip prob.\ $0.15$/step \\
Passive  & $40 \times 50$   & 2{,}000  & $\pm 0.05$ rad & $u \equiv 0$ \\
Eq.\ self-loops (injected) & --- & 200 & $0$ & $u \equiv 0$ \\
\bottomrule
\end{tabular}
\caption{Composition of the training dataset. Total $\approx 27{,}700$
transitions, split 80/10/10 (train/val/test) by \emph{whole episodes}
(no episode straddles two splits).}
\end{table}

Each segment serves a distinct purpose:
\begin{itemize}[leftmargin=2em]
  \item \textbf{Random:} broad state-space coverage so the encoder generalizes;
  \item \textbf{LQR:} near-equilibrium dynamics --- required for $A_{\text{jac}}$
        and $B_{\text{jac}}$ (the linearization at $\zstar$) to be meaningful;
  \item \textbf{Eq-LQR:} very tight initialization teaches
        $f(\zstar, 0) \approx \zstar$ (the fixed-point property);
  \item \textbf{PRBS:} persistent excitation near equilibrium --- needed to
        identify the action-response direction $B$ reliably (LQR uses very
        small $u$, which under-excites $B$);
  \item \textbf{Passive:} natural divergence from near-equilibrium teaches the
        predictor that $\rho(A) > 1$ (the equilibrium is unstable);
  \item \textbf{Equilibrium self-loops:} 200 identical
        $(\text{obs}_{\text{eq}},\ u{=}0,\ \text{obs}_{\text{eq}})$ transitions
        appended directly to the training set, reinforcing the fixed-point via
        \loss{pred} batches.
\end{itemize}

\subsection{TrajectoryDataset}

Returns windows of $H{+}1 = 11$ consecutive frames \emph{from the same episode}
(episode IDs prevent windows from crossing episode boundaries). Frame stacking
at window position $k$:
\begin{equation}
\text{obs\_seq}[k] = \big[\,\text{obs}_{k-1},\ \text{obs}_k\,\big] \in \R^{6 \times 64 \times 64},
\end{equation}
with $\text{obs}_{-1} := \text{obs}_0$ at the start of a window (duplicated frame,
zero-velocity context).

\section{Loss Functions}

Each training batch provides
$\text{obs\_seq} \in \R^{B \times 11 \times 6 \times 64 \times 64}$,
$\text{actions} \in \R^{B \times 10 \times 1}$,
$\text{states} \in \R^{B \times 11 \times 4}$.

\subsection{Forward pass / target construction}

\begin{align}
z_0 &= \text{encoder}_{\text{online}}(\text{obs\_seq}[:,0]) && \text{(gradients flow)} \\
z_{1}, \ldots, z_{10} &= \text{encoder}_{\text{target}}(\text{obs\_seq}[:,1{:}]) && \text{(no\_grad, detached)} \\
z_{\text{targets}} &= [z_1, \ldots, z_{10}].\text{detach}()
\end{align}

\subsection{Predictor unrolling (open-loop)}

Window buffers are initialized as
$z_{\text{buf}} = [z_0, z_0, z_0]$, $\;u_{\text{buf}} = [0, 0]$.
For each step $k = 0, \ldots, 9$:
\begin{align}
\hat z_{k+1} &= f_\theta\big( z_{\text{buf}}[-3{:}],\ u_{\text{buf}}[-3{:}] \big), \\
z_{\text{buf}} &\leftarrow z_{\text{buf}} \,\Vert\, [\hat z_{k+1}].
\end{align}

\begin{quote}
\textcolor{red}{\textbf{Important:}} after step $k=0$ the window contains
\emph{predicted} latents $\hat z$, not re-encoded ground-truth latents. This is
\emph{open-loop} (free-running) unrolling --- prediction errors compound across
the 10-step horizon, and the gradient signal from later steps is dominated by
compounded error rather than clean dynamics information.
\end{quote}

\subsection{$\loss{pred}$ --- Prediction loss \quad ($\lambda = 1.0$)}
\begin{equation}
\loss{pred} = \frac{1}{H} \sum_{k=0}^{H-1} \big\| \hat z_{k+1} - z_{k+1} \big\|^2 .
\end{equation}
The core JEPA objective: the predictor must model latent dynamics such that
its open-loop rollout matches target-encoder embeddings of the true future
frames.

\subsection{$\loss{inv}$ --- Inverse dynamics \quad ($\lambda = 1.0$)}
\begin{equation}
\loss{inv} = \left\| \psi(z_0, z_H) - \frac{\bar u}{10} \right\|^2,
\qquad \bar u = \frac{1}{H}\sum_{k=0}^{H-1} u_k,
\end{equation}
where $\psi : \R^{16} \to \R$ is a 2-layer MLP
($\text{Linear}(16{\to}8) \to \text{ReLU} \to \text{Linear}(8{\to}1)$) and the
constant $10$ is \texttt{inv\_action\_scale} (normalizing the action range
$[-10,10]$). Using the full $H$-step gap ($z_0$ vs.\ $z_H$) instead of
consecutive pairs prevents a trivial "copy" shortcut.

\subsection{$\loss{fp}$ --- Fixed-point loss \quad ($\lambda = 30.0$)}
\begin{equation}
\loss{fp} = \big\| f_\theta\big([\zstar, \zstar, \zstar],\, [0,0,0]\big) - \zstar \big\|^2,
\qquad \zstar = \text{encoder}(\text{obs}_{\text{eq}}).\text{detach}().
\end{equation}
Forces the predictor to treat the upright-equilibrium latent as a fixed point
under zero action. \textbf{This carries the single largest weight in the
system} ($\lambda_{\text{fp}} = 30$, six times larger than $\lambda_{\text{spec}}$
and thirty times larger than $\lambda_{\text{pred}}$) — it is effectively
treated as a near-hard constraint, because without it the predictor could
satisfy $\loss{pred}$ with arbitrary local dynamics near $\zstar$, making the
Jacobian $A_{\text{jac}}$ (linearization at $\zstar$, used for all downstream
control) unreliable.

\subsection{$\loss{spec}$ --- Spectral eigenvalue matching \quad ($\lambda = 5.0$)}

Requires the Jacobian of the predictor at $\zstar$, recomputed every training
step (\texttt{jacobian\_every=1}) via autograd. For a windowed predictor
($W=3$, $d=8$), the augmented \emph{companion-form} Jacobian is
$A_{\text{aug}} \in \R^{24 \times 24}$, $B_{\text{aug}} \in \R^{24 \times 1}$,
constructed as
\begin{equation}
A_{\text{aug}} =
\begin{bmatrix}
0      & I      & 0      & \cdots \\
0      & 0      & I      & \cdots \\
\vdots &        &        & \ddots \\
J_0    & J_1    & \cdots & J_{W-1}
\end{bmatrix},
\qquad
J_k = \frac{\partial \hat z_{t+1}}{\partial z_{\text{window}[k]}} \bigg|_{\zstar}.
\end{equation}
The current-state $(d \times d)$ block used for control is
$A_{\text{jac}} = A_{\text{aug}}[(W{-}1)d{:},\ (W{-}1)d{:}]$.

The loss matches the spectrum of $A_{\text{aug}}$ (all 24 eigenvalues) to the
ground-truth unstable eigenvalues $\Lambda_{\text{unstable}}$ of the
linearized physical cartpole (computed analytically by
\texttt{CartpoleGroundTruth}):
\begin{equation}
\loss{spec} = \sum_{\lambda^\dagger \in \Lambda_{\text{unstable}}}
\ \min_{\lambda_i \in \text{eig}(A_{\text{aug}})} |\lambda_i - \lambda^\dagger|^2
\;+\; 2\big(\rho(A_{\text{aug}}) - \rho_{\text{GT}}\big)^2,
\end{equation}
where $\rho(\cdot)$ denotes spectral radius. The second term additionally
penalizes deviation of the overall growth rate from the ground-truth value
($\rho_{\text{GT}} \approx 1.18$, from $\sqrt{g/\ell}$ discretized at $dt=0.02$s).

\subsection{$\loss{PBH}$ --- Popov--Belevitch--Hautus alignment \quad ($\lambda = 1.0$)}

The PBH controllability test states that an unstable mode (eigenvector $v_u$ of
$A$ with $|\lambda| > 1$) is controllable if and only if $B$ is not orthogonal
to $v_u$. The loss directly enforces this:
\begin{equation}
\loss{PBH} = 1 - \max_{v_u \in V_{\text{unstable}}(A_{\text{aug}})} \cos^2\!\big(\hat b_{\text{eff}},\, v_u\big),
\end{equation}
where $\hat b_{\text{eff}}$ is the unit vector of
$B_{\text{eff,aug}} = B_{\text{aug}} \cdot W_{\text{enc}} \in \R^{24}$
(chaining $\partial z_{t+1}/\partial c_t$ with $\partial c_t/\partial u_t = W_{\text{enc}}$),
and $V_{\text{unstable}}$ is the set of (real parts of) eigenvectors with
$|\lambda| \geq 0.95\,\rho(A_{\text{aug}})$. The loss is $0$ when $B_{\text{eff}}$
perfectly aligns with an unstable mode and $1$ when orthogonal to all of them.

\subsection{$\loss{dynSIG}$ --- Dynamics-aware SIGreg \quad ($\lambda = 2.0$)}

Shapes the encoder's output covariance $\Sigma_z$ toward the
\emph{controllability Gramian} of the local linearization:
\begin{equation}
W_T = \sum_{k=0}^{T_g - 1} A_{\text{jac}}^k\, B_{\text{eff}}\, B_{\text{eff}}^\top\, (A_{\text{jac}}^\top)^k,
\qquad T_g = 5,
\end{equation}
\begin{equation}
\Sigma_{\text{target}} = (1 - \alpha)\, I + \alpha\, \frac{d \cdot W_T}{\operatorname{tr}(W_T)},
\qquad \alpha = 0.3,
\end{equation}
updated by EMA with $\beta = 0.95$ each time the Jacobian is recomputed. The
loss itself is the Epps--Pulley characteristic-function (CF) test (SIGreg, from
LeJEPA 2025) comparing the empirical distribution of random 1-D projections of
$z_0$ against $\mathcal{N}(0, \Sigma_{\text{target}})$, using
\texttt{sigreg\_num\_slices = 32} random projection directions and
\texttt{sigreg\_num\_points = 9} CF evaluation points:
\begin{equation}
\loss{dynSIG} = \frac{1}{S}\sum_{s=1}^{S}
\int \big| \hat\varphi_s(t) - \varphi_{\mathcal{N}(0,\sigma_s^2)}(t) \big|^2 \, w(t)\, dt,
\end{equation}
where $\hat\varphi_s$ is the empirical CF of the $s$-th projection
$\langle z_0, \xi_s \rangle$ and $\sigma_s^2 = \xi_s^\top \Sigma_{\text{target}} \xi_s$.

\textbf{Purpose:} the Gramian encodes which latent directions are reachable by
the control input. Shaping $\Sigma_z$ toward $W_T$ biases the encoder to spread
representational variance into the \emph{controllable} subspace, which should
make downstream LQR/MPC more effective.

\subsection{Total objective}
\begin{equation}
\boxed{
\mathcal{L} =
\underbrace{1.0}_{\lambda_{\text{pred}}} \loss{pred}
+ \underbrace{1.0}_{\lambda_{\text{inv}}} \loss{inv}
+ \underbrace{30.0}_{\lambda_{\text{fp}}} \loss{fp}
+ \underbrace{5.0}_{\lambda_{\text{spec}}} \loss{spec}
+ \underbrace{1.0}_{\lambda_{\text{PBH}}} \loss{PBH}
+ \underbrace{2.0}_{\lambda_{\text{dynSIG}}} \loss{dynSIG}
}
\end{equation}

All other available losses
($\loss{state}$, $\loss{mirror}$, $\loss{anchor}$, $\loss{enc\_anchor}$,
$\loss{sigreg}$, $\loss{varfloor}$, $\loss{temp}$, $\loss{spec\_eig}$,
$\loss{local}$, $\loss{unstable}$) are set to $\lambda = 0$ in
\texttt{cartpole\_v2\_fullspec.yaml} and do not contribute.

\subsection{Optimizer}

Adam, with parameter groups:
\begin{itemize}[leftmargin=2em]
  \item encoder: $\text{lr} = 10^{-4}$, weight decay $= 10^{-4}$;
  \item action encoder: $\text{lr} = 10^{-4}$, weight decay $= 0$;
  \item predictor: $\text{lr} = 10^{-4}$ (\texttt{predictor\_lr\_mult=1.0}), weight decay $= 10^{-4}$;
\end{itemize}
with cosine-annealing schedule, $T_{\max} = 200$ epochs, $\eta_{\min} = 10^{-5}$.

\section{Potential Failure Modes}

Given the precise loss landscape above, the most likely sources of poor
downstream control performance are:

\begin{enumerate}[leftmargin=2em]
  \item \textbf{Open-loop unrolling in $\loss{pred}$.} The predicted
        $\hat z$ (not the re-encoded ground truth) fills the window buffer for
        $k \geq 1$. Errors compound multiplicatively over the 10-step horizon;
        gradient contributions from later steps are dominated by compounded
        error rather than clean one-step dynamics signal.

  \item \textbf{$z^*$ mismatch between training and evaluation.} At training,
        $\text{obs}_{\text{eq},2f} = [\text{obs}_{\text{eq}}, \text{obs}_{\text{eq}}]$
        (duplicated frame). At evaluation, the very first frame after
        \texttt{reset\_to\_state(0)} may differ subtly (e.g.\ rendering jitter,
        physics solver state). Any drift in the encoded $\zstar$ propagates
        directly into $A_{\text{jac}}$, $B_{\text{jac}}$, and hence the LQR
        gain.

  \item \textbf{$\loss{spec}$ targets all 24 eigenvalues of $A_{\text{aug}}$,}
        but the ground truth only specifies the \emph{physical} unstable
        eigenvalue ($\lambda^\dagger \approx 1.18$). The remaining $\sim$23
        "history" eigenvalues (artifacts of the companion-form padding
        structure for $W=3$) receive gradient pressure toward the nearest GT
        eigenvalue with no principled target, potentially distorting the
        Jacobian's structure in directions irrelevant to control.

  \item \textbf{$\lambda_{\text{fp}} = 30$ dominates the objective.} At 30$\times$
        the weight of $\loss{pred}$, the predictor is primarily optimizing to
        be a fixed point at $\zstar$ rather than to model dynamics broadly. It
        may learn a narrow "gravity well" around $\zstar$ that satisfies
        $\loss{fp}$ and the near-equilibrium portion of $\loss{pred}$, without
        generalizing to states further away --- exactly where the
        cart must operate during a swing-up or large-disturbance recovery.

  \item \textbf{Companion-form Jacobian dimensionality.} With $W=3$, $d=8$,
        the augmented system is $24$-dimensional, but the physical system is
        $4$-dimensional (and the latent system is nominally $8$-dimensional,
        $d = 2 \times n_{\text{state}}$ for frame-stacked position+velocity
        encoding). The gap between "true" and "augmented" dimensionality means
        most of $A_{\text{aug}}$'s spectral content has no physical
        counterpart, yet $\loss{spec}$ and $\loss{PBH}$ both operate on the
        full augmented matrix.
\end{enumerate}

\end{document}
